\chapter{Dictionaries via Hashing}

\tikzset{
  sl/.style={draw,minimum width=9mm,minimum height=6mm,font=\small},
  ix/.style={font=\scriptsize,gray},
  ch/.style={draw,rounded corners,fill=blue!8,minimum height=5mm,font=\scriptsize,inner sep=2pt}
}

\section{Direct Addressing}

If keys are integers in $\{0, \ldots, M-1\}$, store the dictionary in an array
$A$ of size $M$, with $A[k]$ holding the value for key $k$ (or \code{nil}).
\code{search}, \code{insert}, \code{delete} are all $\Theta(1)$, but space is
$\Theta(M)$, which is wasteful when $n \ll M$ and impossible when keys come from a
huge universe $U$.

\section{Hashing}

\begin{definition}
A \emph{hash function} $h : U \to \{0, \ldots, M-1\}$ maps keys to slots of a
\emph{hash table} $T$ of size $M$. Ideally key $k$ is stored at $T[h(k)]$. A
\emph{collision} occurs when $h(k_1) = h(k_2)$ for $k_1 \ne k_2$.
\end{definition}

Since $|U| \gg M$, collisions are unavoidable. Two decisions: choose a good $h$,
and choose a collision-resolution strategy.

\begin{definition}
The \emph{load factor} is $\alpha = n / M$.
\end{definition}

We keep $\alpha$ bounded by \emph{rehashing}: when $\alpha$ exceeds a threshold,
allocate a table of roughly twice the size, choose a new hash function, and
reinsert everything ($\Theta(n + M)$). Rehashing happens rarely enough that the
amortized cost per operation is $\bigO(1)$. Similarly shrink when $\alpha$ gets too
small, so space stays $\Theta(n)$.

\paragraph{Uniform hashing assumption.} For analysis, assume each key is
independently equally likely to hash to any slot. This is not literally true
for a fixed $h$, which is why we later discuss universal hashing.

\section{Separate Chaining}

Each slot holds a bucket (usually an unsorted linked list) of all keys hashing to it.
\begin{itemize}
  \item \code{search(k)}: scan bucket $T[h(k)]$.
  \item \code{insert(k,v)}: add to the front of bucket $T[h(k)]$: $\Theta(1)$.
  \item \code{delete(k)}: search, then unlink.
\end{itemize}
Under uniform hashing, the expected bucket length is $\alpha$, so the expected
cost of \code{search}/\code{delete} is $\Theta(1 + \alpha)$. With $\alpha \in
\Theta(1)$ (maintained by rehashing), all operations are $\Theta(1)$ expected,
while the worst case is $\Theta(n)$ (all keys in one bucket).

With the keys of the linear-probing example below ($M = 11$,
$h(k) = k \bmod 11$), chaining gives:
\begin{center}
\begin{tikzpicture}
  \node[draw,minimum width=6mm,minimum height=5mm] (t0) at (0,0.0) {};
  \node[ix,anchor=east] at (-0.4,0.0) {0};
  \node[ch] (c0_0) at (1.2,0.0) {44};
  \draw[-{Stealth}] (t0.center) -- (c0_0);
  \node[ch] (c0_1) at (2.3,0.0) {22};
  \draw[-{Stealth}] (c0_0) -- (c0_1);
  \fill (t0.center) circle (1.2pt);
  \node[draw,minimum width=6mm,minimum height=5mm] (t1) at (0,-0.5) {};
  \node[ix,anchor=east] at (-0.4,-0.5) {1};
  \node[draw,minimum width=6mm,minimum height=5mm] (t2) at (0,-1.0) {};
  \node[ix,anchor=east] at (-0.4,-1.0) {2};
  \node[draw,minimum width=6mm,minimum height=5mm] (t3) at (0,-1.5) {};
  \node[ix,anchor=east] at (-0.4,-1.5) {3};
  \node[draw,minimum width=6mm,minimum height=5mm] (t4) at (0,-2.0) {};
  \node[ix,anchor=east] at (-0.4,-2.0) {4};
  \node[ch] (c4_0) at (1.2,-2.0) {59};
  \draw[-{Stealth}] (t4.center) -- (c4_0);
  \fill (t4.center) circle (1.2pt);
  \node[draw,minimum width=6mm,minimum height=5mm] (t5) at (0,-2.5) {};
  \node[ix,anchor=east] at (-0.4,-2.5) {5};
  \node[draw,minimum width=6mm,minimum height=5mm] (t6) at (0,-3.0) {};
  \node[ix,anchor=east] at (-0.4,-3.0) {6};
  \node[draw,minimum width=6mm,minimum height=5mm] (t7) at (0,-3.5) {};
  \node[ix,anchor=east] at (-0.4,-3.5) {7};
  \node[ch] (c7_0) at (1.2,-3.5) {73};
  \draw[-{Stealth}] (t7.center) -- (c7_0);
  \fill (t7.center) circle (1.2pt);
  \node[draw,minimum width=6mm,minimum height=5mm] (t8) at (0,-4.0) {};
  \node[ix,anchor=east] at (-0.4,-4.0) {8};
  \node[ch] (c8_0) at (1.2,-4.0) {41};
  \draw[-{Stealth}] (t8.center) -- (c8_0);
  \fill (t8.center) circle (1.2pt);
  \node[draw,minimum width=6mm,minimum height=5mm] (t9) at (0,-4.5) {};
  \node[ix,anchor=east] at (-0.4,-4.5) {9};
  \node[ch] (c9_0) at (1.2,-4.5) {31};
  \draw[-{Stealth}] (t9.center) -- (c9_0);
  \fill (t9.center) circle (1.2pt);
  \node[draw,minimum width=6mm,minimum height=5mm] (t10) at (0,-5.0) {};
  \node[ix,anchor=east] at (-0.4,-5.0) {10};
  \node[ch] (c10_0) at (1.2,-5.0) {32};
  \draw[-{Stealth}] (t10.center) -- (c10_0);
  \fill (t10.center) circle (1.2pt);
  \node[font=\scriptsize,anchor=west,align=left] at (3.6,-0.25) {22 and 44 both hash to 0;\\44 was inserted later, at the front};
\end{tikzpicture}
\end{center}

\section{Open Addressing}

Store all items directly in the table, at most one per slot ($\alpha \le 1$).
Each key has a \emph{probe sequence} $h(k, 0), h(k, 1), \ldots$ of slots to try.
\begin{lstlisting}
probe-insert(T, (k, v)):
  for j <- 0 to M-1 do
    if T[h(k, j)] is empty or "deleted" then
      T[h(k, j)] <- (k, v); return "success"
  return "failure to insert"       // need to rehash

probe-search(T, k):
  for j <- 0 to M-1 do
    if T[h(k, j)] is empty then return "not found"
    else if T[h(k, j)] has key k then return T[h(k, j)]
    // else: occupied by another key or "deleted": keep probing
  return "not found"
\end{lstlisting}

\paragraph{Deletion.} We cannot simply empty a slot, since that would break
probe sequences of other keys passing through it. Instead use \emph{lazy
deletion}: mark the slot as ``deleted''. Search continues past such slots;
insert may reuse them. Rehash when too many deleted markers accumulate.

\subsection{Linear Probing}

$h(k, j) = (h(k) + j) \bmod M$. Simple and cache-friendly, but suffers from
\emph{primary clustering}: occupied runs grow, and a key hashing anywhere into a
run extends it, so long runs get longer.

\begin{example}
$M = 11$, $h(k) = k \bmod 11$. Insert 41, 22, 44, 59, 32, 31, 73.
\begin{itemize}
  \item $41 \to 8$; $22 \to 0$; $44 \to 0$ collides, goes to 1;
  \item $59 \to 4$; $32 \to 10$; $31 \to 9$;
  \item $73 \to 7$.
\end{itemize}
Now inserting 19 ($h = 8$) probes 8, 9, 10, 0, 1 and lands in 2: the clusters at
7--10 and 0--1 merged into one.
\begin{center}
\begin{tikzpicture}
  \node[sl,fill=blue!8] (s0) at (0.0,0) {22};
  \node[ix] at (0.0,-0.55) {0};
  \node[sl,fill=blue!8] (s1) at (0.9,0) {44};
  \node[ix] at (0.9,-0.55) {1};
  \node[sl,fill=orange!25] (s2) at (1.8,0) {19};
  \node[ix] at (1.8,-0.55) {2};
  \node[sl,fill=white] (s3) at (2.7,0) {};
  \node[ix] at (2.7,-0.55) {3};
  \node[sl,fill=blue!8] (s4) at (3.6,0) {59};
  \node[ix] at (3.6,-0.55) {4};
  \node[sl,fill=white] (s5) at (4.5,0) {};
  \node[ix] at (4.5,-0.55) {5};
  \node[sl,fill=white] (s6) at (5.4,0) {};
  \node[ix] at (5.4,-0.55) {6};
  \node[sl,fill=blue!8] (s7) at (6.3,0) {73};
  \node[ix] at (6.3,-0.55) {7};
  \node[sl,fill=blue!8] (s8) at (7.2,0) {41};
  \node[ix] at (7.2,-0.55) {8};
  \node[sl,fill=blue!8] (s9) at (8.1,0) {31};
  \node[ix] at (8.1,-0.55) {9};
  \node[sl,fill=blue!8] (s10) at (9.0,0) {32};
  \node[ix] at (9.0,-0.55) {10};
  \draw[-{Stealth},red!70!black,thick] (s8.north) to[out=60,in=120] (s9.north);
  \draw[-{Stealth},red!70!black,thick] (s9.north) to[out=60,in=120] (s10.north);
  \draw[-{Stealth},red!70!black,thick] (s10.north) to[out=90,in=90,looseness=0.4] (s0.north);
  \draw[-{Stealth},red!70!black,thick] (s0.north) to[out=60,in=120] (s1.north);
  \draw[-{Stealth},red!70!black,thick] (s1.north) to[out=60,in=120] (s2.north);
  \node[font=\scriptsize,red!70!black] at (4.5,1.55) {probes for 19: 8, 9, 10, 0, 1, then 2 is free};
\end{tikzpicture}
\end{center}
\end{example}

\subsection{Quadratic Probing}

$h(k,j) = (h(k) + c_1 j + c_2 j^2) \bmod M$. Avoids primary clustering but has
\emph{secondary clustering} (keys with the same $h(k)$ share a probe sequence),
and must be designed carefully so the sequence visits enough slots (e.g.\ $M$
prime and $\alpha < 1/2$ guarantees an empty slot is found).

\subsection{Double Hashing}

Use two independent hash functions:
\[
h(k, j) = (h_0(k) + j\, h_1(k)) \bmod M,
\]
where $h_1(k) \neq 0$ and is relatively prime to $M$ (e.g.\ $M$ prime) so the
probe sequence is a permutation of all slots. Keys that collide at $h_0$ usually
get different step sizes, avoiding both kinds of clustering. A common choice is
the multiplicative method for $h_1$: $h_1(k) = \floor{M' (kA - \floor{kA})} + 1$
with $A \approx (\sqrt5-1)/2$.

\subsection{Cuckoo Hashing}

Use two hash functions $h_0, h_1$ and two tables $T_0, T_1$ (or one table).
Every key $k$ is stored in either $T_0[h_0(k)]$ or $T_1[h_1(k)]$ --- nowhere else.
\begin{itemize}
  \item \code{search} and \code{delete} check exactly two slots: $\Theta(1)$
        \emph{worst case}.
  \item \code{insert} places $k$ in $T_0[h_0(k)]$; if that slot was occupied by
        $k'$, then $k'$ is kicked out and moves to its alternative slot in the
        other table, possibly kicking out another key, and so on.
\end{itemize}
\begin{lstlisting}
cuckoo-insert(k, v):
  i <- 0
  do at most 2n times:
    if T_i[h_i(k)] is empty then
      T_i[h_i(k)] <- (k, v); return "success"
    swap (k, v) with T_i[h_i(k)]     // evict the occupant
    i <- 1 - i                       // evicted key goes to the other table
  return "failure"                   // a cycle: rehash with new h0, h1
\end{lstlisting}
If insert loops too long, there is a cycle and we rehash with new hash functions
(and possibly a bigger table). With $\alpha < 1/2$ (across both tables) and
random-like hash functions, the expected insertion time is $\bigO(1)$.

\begin{example}
$M = 11$, $h_0(k) = k \bmod 11$, $h_1(k) = \floor{11(0.618k - \floor{0.618k})}$.
Insert 51: $h_0 = 7$, goes to $T_0[7]$. Insert 95: $h_0 = 7$ is taken, so 95
takes $T_0[7]$ and 51 is evicted to $T_1[h_1(51)]$. Here
$h_1(51) = \floor{11 \cdot 0.518} = 5$:
\begin{center}
\begin{tikzpicture}
  \node[sl,fill=white] (u0_0) at (0.0,0.0) {};
  \node[ix] at (0.0,0.5) {0};
  \node[sl,fill=white] (u0_1) at (0.9,0.0) {};
  \node[ix] at (0.9,0.5) {1};
  \node[sl,fill=white] (u0_2) at (1.8,0.0) {};
  \node[ix] at (1.8,0.5) {2};
  \node[sl,fill=white] (u0_3) at (2.7,0.0) {};
  \node[ix] at (2.7,0.5) {3};
  \node[sl,fill=white] (u0_4) at (3.6,0.0) {};
  \node[ix] at (3.6,0.5) {4};
  \node[sl,fill=white] (u0_5) at (4.5,0.0) {};
  \node[ix] at (4.5,0.5) {5};
  \node[sl,fill=white] (u0_6) at (5.4,0.0) {};
  \node[ix] at (5.4,0.5) {6};
  \node[sl,fill=blue!8] (u0_7) at (6.3,0.0) {95};
  \node[ix] at (6.3,0.5) {7};
  \node[sl,fill=white] (u0_8) at (7.2,0.0) {};
  \node[ix] at (7.2,0.5) {8};
  \node[sl,fill=white] (u0_9) at (8.1,0.0) {};
  \node[ix] at (8.1,0.5) {9};
  \node[sl,fill=white] (u0_10) at (9.0,0.0) {};
  \node[ix] at (9.0,0.5) {10};
  \node[font=\small,anchor=east] at (-0.6,0.0) {$T_0$};
  \node[sl,fill=white] (u1_0) at (0.0,-1.6) {};
  \node[ix] at (0.0,-2.1) {0};
  \node[sl,fill=white] (u1_1) at (0.9,-1.6) {};
  \node[ix] at (0.9,-2.1) {1};
  \node[sl,fill=white] (u1_2) at (1.8,-1.6) {};
  \node[ix] at (1.8,-2.1) {2};
  \node[sl,fill=white] (u1_3) at (2.7,-1.6) {};
  \node[ix] at (2.7,-2.1) {3};
  \node[sl,fill=white] (u1_4) at (3.6,-1.6) {};
  \node[ix] at (3.6,-2.1) {4};
  \node[sl,fill=blue!8] (u1_5) at (4.5,-1.6) {51};
  \node[ix] at (4.5,-2.1) {5};
  \node[sl,fill=white] (u1_6) at (5.4,-1.6) {};
  \node[ix] at (5.4,-2.1) {6};
  \node[sl,fill=white] (u1_7) at (6.3,-1.6) {};
  \node[ix] at (6.3,-2.1) {7};
  \node[sl,fill=white] (u1_8) at (7.2,-1.6) {};
  \node[ix] at (7.2,-2.1) {8};
  \node[sl,fill=white] (u1_9) at (8.1,-1.6) {};
  \node[ix] at (8.1,-2.1) {9};
  \node[sl,fill=white] (u1_10) at (9.0,-1.6) {};
  \node[ix] at (9.0,-2.1) {10};
  \node[font=\small,anchor=east] at (-0.6,-1.6) {$T_1$};
  \draw[-{Stealth},thick,red!70!black] (u0_7.south) -- node[right=3pt,font=\scriptsize]{51 evicted} (u1_5.north);
\end{tikzpicture}
\end{center}
\end{example}

\subsection{Expected Probes}

Under uniform hashing, for open addressing with load factor $\alpha < 1$:
\begin{center}
\begin{tabular}{lll}
\toprule
Strategy & Unsuccessful search / insert & Successful search\\
\midrule
Linear probing & $\frac{1}{(1-\alpha)^2}$ & $\frac{1}{1-\alpha}$\\
Double hashing (ideal) & $\frac{1}{1-\alpha}$ & $\frac{1}{\alpha}\ln\frac{1}{1-\alpha}$\\
Cuckoo hashing & $\frac{\alpha}{(1-2\alpha)^2}$ (insert) & $1$ (worst case $\le 2$)\\
Separate chaining & $1 + \alpha$ & $1 + \alpha/2$\\
\bottomrule
\end{tabular}
\end{center}
(Up to constants and lower-order terms.) The key point: all are $\bigO(1)$ when
$\alpha$ is bounded away from the critical value (1 for probing, $1/2$ for cuckoo).

\begin{center}
\begin{tikzpicture}[x=8cm,y=0.4cm]
  \begin{scope}
  \clip (0,0) rectangle (1,12.5);
  \draw[very thick,red!70!black,domain=0:0.75,samples=60] plot (\x,{1/((1-\x)*(1-\x))});
  \draw[very thick,blue!60!black,domain=0:0.93,samples=60] plot (\x,{1/(1-\x)});
  \draw[very thick,green!50!black,domain=0:1,samples=2] plot (\x,{1+\x});
  \end{scope}
  \draw[-{Stealth}] (0,0) -- (1.05,0) node[right,font=\scriptsize]{$\alpha$};
  \draw[-{Stealth}] (0,0) -- (0,13) node[above,font=\scriptsize]{expected probes};
  \foreach \t in {0.25,0.5,0.75,1} \draw (\t,0) -- (\t,-0.25) node[below,font=\scriptsize]{\t};
  \foreach \t in {4,8,12} \draw (0,\t) -- (-0.01,\t) node[left,font=\scriptsize]{\t};
  \node[font=\scriptsize,red!70!black,anchor=east] at (0.7,10) {linear probing $\frac{1}{(1-\alpha)^2}$};
  \node[font=\scriptsize,blue!60!black,anchor=west] at (0.93,11) {double hashing $\frac{1}{1-\alpha}$};
  \node[font=\scriptsize,green!50!black,anchor=west] at (0.75,2.6) {chaining $1+\alpha$};
\end{tikzpicture}
\end{center}
Unsuccessful search: all curves are flat for small $\alpha$, but the probing curves blow up as
$\alpha \to 1$.

\begin{proof}[Sketch for ideal double hashing]
If each probe is to a uniformly random slot not yet probed, each probe finds an
empty slot with probability at least $1 - \alpha$. The number of probes is then
dominated by a geometric random variable with mean $\frac1{1-\alpha}$.
\end{proof}

\section{Choosing a Hash Function}

Goals: fast to compute and spreads keys uniformly, even for structured
(non-random) keys.
\begin{itemize}
  \item \textbf{Modular method:} $h(k) = k \bmod M$ with $M$ prime, not close to
        a power of 2 (otherwise $h$ depends only on low-order bits).
  \item \textbf{Multiplicative method:} $h(k) = \floor{M (kA - \floor{kA})}$
        for a constant $0 < A < 1$; Knuth suggests $A = \varphi - 1 \approx 0.618$.
        Here $M$ can be any size.
  \item \textbf{Strings:} interpret $s = s_0 s_1 \ldots s_{p-1}$ as a number in
        base $R$ (e.g.\ 255) and reduce mod $M$, using Horner's rule:
        $h(s) = (\cdots((s_0 R + s_1) R + s_2)\cdots) \bmod M$, reducing at each step.
\end{itemize}

\subsection{Universal Hashing}

For any fixed $h$, an adversary can choose $n$ keys that all collide, forcing
$\Theta(n)$ per operation. The fix is randomization: choose $h$ at random from a
family when the table is created (and at each rehash).

\begin{definition}
A family $\mathcal H$ of hash functions $U \to \{0,\ldots,M-1\}$ is
\emph{universal} if for every pair of distinct keys $k_1 \ne k_2$,
\[
\Pr_{h \in \mathcal H}[h(k_1) = h(k_2)] \le \frac1M.
\]
\end{definition}

\begin{theorem}
Let $p > |U|$ be prime and $\mathcal H = \{h_{a,b}(k) = ((ak + b) \bmod p) \bmod M
: 1 \le a < p,\ 0 \le b < p\}$. Then $\mathcal H$ is universal.
\end{theorem}

\begin{theorem}
With $h$ drawn from a universal family and separate chaining, the expected
length of the bucket for any key $k$ is at most $1 + \alpha$. Thus every
operation takes $\bigO(1 + \alpha)$ expected time for \emph{every} input.
\end{theorem}
\begin{proof}
For each other key $k'$ in the table, let $X_{k'}$ be 1 if $h(k') = h(k)$.
$\mathbb E[X_{k'}] \le 1/M$. By linearity, the expected number of colliding keys
is $\le n/M = \alpha$; add 1 for $k$ itself.
\end{proof}

\section{Hashing vs.\ Balanced Search Trees}

\begin{center}
\begin{tabular}{lll}
\toprule
 & Hash table & AVL tree\\
\midrule
search/insert/delete & $\bigO(1)$ expected & $\Theta(\log n)$ worst case\\
worst case & $\Theta(n)$ (cuckoo: $\bigO(1)$ search) & $\Theta(\log n)$\\
ordered operations (min, successor, range) & not supported & supported\\
space & wasted empty slots & pointers + heights\\
\bottomrule
\end{tabular}
\end{center}

\section{Summary}
\begin{itemize}
  \item Direct addressing is $\Theta(1)$ but needs space $\Theta(|U|)$.
  \item Hashing maps a large universe to $M$ slots; keep $\alpha = n/M$ bounded by
        rehashing (amortized $\bigO(1)$).
  \item Chaining: expected $\Theta(1 + \alpha)$.
  \item Open addressing: linear probing (primary clustering), quadratic probing
        (secondary clustering), double hashing; lazy deletion.
  \item Cuckoo hashing: $\bigO(1)$ worst-case search and delete; expected $\bigO(1)$
        insert when $\alpha < 1/2$.
  \item Universal hashing gives expected guarantees for every input.
\end{itemize}
